\documentclass[10pt,a4paper,fontset=none]{ctexart}
\setCJKmainfont{Songti SC}
\setCJKsansfont{Heiti SC}
\setCJKmonofont{Songti SC}

\usepackage{amsmath,amssymb,amsthm,geometry}
\usepackage{booktabs}
\usepackage[hidelinks]{hyperref}

\geometry{margin=2cm}

\title{标题}
\author{LeyuDame}
\date{\today}

\newtheorem{theorem}{定理}
\newtheorem{lemma}{引理}
\newtheorem{proposition}{命题}
\newtheorem{example}{例}
\newtheorem{remark}{注}

\begin{document}

\maketitle

\begin{abstract}
这里用 3--6 句话说明本文讨论的问题、核心方法和主要收获。
\end{abstract}

\section{引言}

从一个自然问题、典型题目或教学困惑切入，说明本文要解决什么问题。

本文结构如下：第 2 节回顾必要的预备知识；第 3 节给出核心方法；第 4 节讨论典型应用；最后进行总结。

\section{预备知识}

这里只放后文真正需要使用的定义、公式、定理和记号。

\section{核心方法}

\subsection{关键观察}

说明问题的关键结构。

\subsection{方法推导}

给出主要推导过程。

\subsection{主要结论}

提炼出可以复用的结论或技巧。

\section{应用}

\subsection{基础应用}

\textbf{题目.}

题目内容。

\textbf{分析.}

说明为什么想到本文的核心方法。

\textbf{解答.}

写出完整解答。

\textbf{点评.}

总结本题体现的触发条件或方法迁移。

\subsection{综合应用}

\textbf{题目.}

题目内容。

\textbf{分析.}

说明关键观察。

\textbf{解答.}

写出完整解答。

\textbf{点评.}

指出本题与前一题的联系和变化。

\section{总结}

总结本文的核心方法、适用场景和识别方式。

\section*{参考资料}

\begin{enumerate}
  \item 作者，资料名称，出版信息或网站信息，年份。
  \item 链接：\url{https://example.com}
\end{enumerate}

\end{document}
